\section*{编写旨趣}

在本书中，我们将以一种模型无关的方式，借助一个为多种模型所共同满足的公共
公理框架，从第一性原理出发发展\trInfinityCategory{}的理论。与先前对\trInfinityCategory{}理论的
「分析式」（analytic）处理——其核心\trCategory{}概念要参照某个特定模型的坐标来
定义——相对照，我们的取径是「综合式」（synthetic）的：从可以在众多模型中
同时诠释的定义出发，而我们的证明随之对这些模型一体适用。虽为综合式，我们的
工作却不流于图式化或含糊其辞——不会把如何将一切彻底精确化的细节留给
「专家们」，也不会一路下去尽是乌龟驮着。\footnote{一种不够严格的「模型
无关」式∞-范畴论表述，会遭遇无穷倒退的问题：无穷维\trCategory{}本身就是某个外围
无穷维\trCategory{}的对象，而发展前者的理论时，人们难免想动用后者的理论。我们选用
一个非常具体的模型，来充当那个由\trInfinityCategory{}构成、在实践中频繁出现、并为便利
相对简单的证明而设计的外围 $(\infty,2)$-\trCategory{}，从而避开这个问题。
$(\infty,2)$-\trCategory{}的理论尚在襁褓之中，我们乐于以这种方式快刀斩断戈尔迪
之结。}\yizhu{「一路下去尽是乌龟」（turtles all the way down）
典出西方「世界驮在龟背上、龟下复有龟」的传说，喻无穷倒退、层层无尽。}
恰恰相反，我们从一份简短而明白列举的公理清单出发证明定理，因而我们的结论
对任何满足这些公理的\trInfinityCategory{}模型都有效。

这个综合式的理论可以在任何\textbf{\trInfinityCosmos}（∞-cosmos）中发展；
\trInfinityCosmos{}将\trInfinityCategory{}栖居其间、充作对象的那个宇宙公理化。为了使定理的陈述能直接
提示其自然诠释，我们把「\trInfinityCategory」重新铸造为一个技术术语：意指某个（通常固定
取定的）\trInfinityCosmos{}中的对象。$(\infty,1)$-\trCategory{}的若干常见模型\footnote{
\trQuasiCategory{}、\trCompleteSegalSpaces{}、\trSegalCategories{}与\trOneComplicialSets{}（即自然带标记的%
\trQuasiCategory）都定义了某个\trInfinityCosmos{}中的\trInfinityCategory{}。}是此种意义上的\trInfinityCategory{}，但我们的%
\trInfinityCategory{}还包括某些 $(\infty,n)$-\trCategory\footnote{$n$-\trQuasiCategory{}、$\Theta_n$-空间、
迭代的\trCompleteSegalSpaces{}与 $n$-复状集同样定义了某个\trInfinityCosmos{}中的\trInfinityCategory{}；饱和
（旧称「弱」）复状集亦然——它是 $(\infty,\infty)$-\trCategory{}的一个模型。}的
模型，以及上述一切的\trFibration{}版本。于是，这些对象个个都是此种意义上的\trInfinityCategory{}，
而我们的定理对它们全体适用。\footnote{不过，在某种意义上，我们的许多定义是
为对象为 $(\infty,1)$-\trCategory{}的那些\trInfinityCosmos{}而优化的。定义 1.2.26 引入的%
\trDiscrete{}\trInfinityCategory{}概念便是一个佳例：在 $(\infty,1)$-\trCategory{}的\trInfinityCosmos{}中，
\trDiscrete\trInfinityCategory{}恰是\trInfinityGroupoid{}。尽管这在 $(\infty,n)$-\trCategory{}的\trInfinityCosmos{}中并不成立，
我们仍将把这一概念用到某些\trExotic\trInfinityCosmos{}中（例如定义 7.4.1）。}「\trInfinityCategory」
一词的这种用法，意在取中于两种既有用法之间：经典的用法泛指弱无穷维\trCategory{}的
任何变种，而通常的用法往往专指\trQuasiCategory{}或\trCompleteSegalSpaces{}。

\trInfinityCategory{}理论的许多发展，并非在全整的\trInfinityCosmos{}中进行，而是在一个我们称为
\textbf{\trHomotopyTwoCategory}（homotopy 2-category）的商结构中进行；
取此名是因为，\trInfinityCosmos{}有点像一个由富集\trModelCategory{}中的\trFibrant{}对象构成的%
\trCategory{}，而\trHomotopyTwoCategory{}则是其\trHomotopyCategory{}（homotopy category）的\trCategory{}化
（categorification）。\trHomotopyTwoCategory{}是一个严格的 2-\trCategory——正如由%
\trCategory{}、\textbf{\trFunctor}（functor）与
\textbf{\trNaturalTransformation}（natural transformation）构成的 2-\trCategory
\footnote{事实上，这又是一个特例：存在一个\trInfinityCosmos{}，其对象是通常的%
\trCategory{}，而其\trHomotopyTwoCategory{}正是那个由\trCategory{}、\trFunctor{}与%
\trNaturalTransformation{}构成的通常\trCategory{}。这个 2-\trCategory{}与范畴论本身一样古老，
见于 Eilenberg 与 Mac Lane 的奠基性论文 [42]。}——正因如此，\trInfinityCategory{}理论中
的奠基性证明与通常范畴论的经典基础极为相似，只是它们所刻画的泛性质
（universal property）略弱于严格 1-\trCategory{}的熟悉情形，例如\trInfinityCategory{}之间的\trFunctor
何时定义一个\textbf{\trCartesianFibration}（cartesian fibration）。

在为\trInfinityCosmos{}遴选公理时，本有许多别的选择余地。我们的指导原则之一——诚然，
这与\trHomotopy{}式高阶范畴论的常见设定颇相径庭——是让自己得以尽可能严格地工作，
以求缩短并简化证明。这些选择的后果是：\trInfinityCosmos{}中的\trInfinityCategory{}，连同它们之间的%
\trFunctor{}与\trNaturalTransformation{}，装配成一个 2-\trCategory{}，而非双\trCategory
（bicategory）。为帮助我们达成这种反直觉的严格性，每个\trInfinityCosmos{}都自带一个
指定的、\trInfinityCategory{}之间的映射类，称为\textbf{\trIsofibration}
（isofibration）。\trIsofibration{}并没有\trHomotopy{}论上的含义——\trInfinityCategory{}之间的任何%
\trFunctor{}都\trEquivalence{}于一个陪域（codomain）相同的\trIsofibration——
然而，\trIsofibration{}使我们得以在\trInfinityCategory{}之间考量严格交换的图，并使我们得以要求
此类图的\trLimit{}在单纯富集\trIsomorphism{}的意义下满足泛性质。这两个特点都并非%
\trInfinityCategory{}理论的发展所必需；类似的证明在更弱的设定下同样行得通，只是要在考量
高维\trCell{}的相干性上多花时间。

在第 I 部分中，我们定义并发展\trInfinityCategory{}之间的\trEquivalence{}与
\textbf{\trAdjunction}（adjunction）、\trInfinityCategory{}中的\trLimit{}与
\textbf{\trColimit}（colimit），以及\trCartesianFibration{}、余\trCartesianFibration
及其离散变体——并为这些对象证明了\trYonedaLemma{}的一个版本。这些结果大多
在\trHomotopyTwoCategory{}的舒适庇护之下发展。在一篇间章（interlude）中，我们岔入抽象%
\trInfinityCosmos{}论，以便对任何\trInfinityCosmos{}中都存在的整类\trLimit{}构造给出更仔细的交代。这一
分析被用来发展更多的\trInfinityCosmos{}实例：其对象为\trPointed\trInfinityCategory{}、或稳定
的\trInfinityCategory{}、或给定\trInfinityCosmos{}中的\trCartesianFibration{}与余笛卡尔%
\trFibration{}。\footnote{性急的读者可以跳过这篇间章，径直相信任何\trInfinityCosmos{}都会衍生
出各种其他的\trInfinityCosmos{}，而不妨碍其对后文的理解——只是会错过一些乐趣。}

这份关于\trInfinityCategory{}范畴论的基本叙述尚缺一角：对\trAdjunction{}每个\trInfinityCategory{}的「hom」双\trFunctor
（bifunctor）的令人满意的处理——它正是我们称为\textbf{\trModule}
（module）之物的原型实例。对新手来说，一道有益的练习是：在某个选定的模型
中，定义那个取值于\trInfinityGroupoid{}的 hom 双\trFunctor{}。而有教益的则在于获知：这个对通常
范畴论而言如此根本的构造，做起来竟难到令人却步。\footnote{拟范畴论的专家
知道该用 Lurie 的拉直–反拉直构造 [78, 2.2.1.2]，或 Cisinski 的泛左\trFibration
[28, 5.2.8]，以及扭箭头\trQuasiCategory{}。}在我们的公理化中，\trInfinityCosmos{}里的任何\trInfinityCategory
都有一个相伴的箭头\trInfinityCategory{}，配备着定义域（domain）投射\trFunctor{}与陪域投射%
\trFunctor{}，二者以彼此相容的方式分别定义\trCartesianFibration{}与余笛卡尔%
\trFibration{}。这样的\trModule——它们自身也装配成一个\trInfinityCosmos——为把泛性质编码为%
\trFibered{}\trEquivalence{}提供了便利的载体。在第 II 部分中，我们将发展\trInfinityCategory{}之间%
\trModule{}的演算，并应用它来定义和研究逐点\trKanExtension{}（pointwise Kan
extensions）。这给了我们一个机会，把第 I 部分中证明的泛性质重新打包，归入%
\trInfinityCategory{}的「形式范畴论」（formal category theory）。

这项工作全部是「模型中立」（model-agnostic）的：对任何特定\trInfinityCosmos{}之规格的
细节视而不见。在第 III 部分中，我们证明\trInfinityCategory{}的范畴论在一种精确的意义下
还是「模型无关的」：一切\trCategory{}概念都被任何「\trChangeOfModel」\trFunctor——即定义了我们所谓
的\textbf{宇宙论\trBiequivalence}（cosmological biequivalence）的\trFunctor
——保持、反射并创造。这个\textbf{\trModelIndependence}（model
independence）定理比公理框架初看所提示的更为强力：它还允许我们把用分析式
技术证明的定理迁移到一切\trBiequivalence{}的\trInfinityCosmos{}。例如，对象为 $(\infty,1)$-\trCategory{}之
模型的那四个\trInfinityCosmos{}全都\trBiequivalence{}。\footnote{一个密切相关的观察是：Joyal 与
Tierney 在 [64] 中构造的\trQuasiCategory{}、\trCompleteSegalSpaces{}与\trSegalCategories{}之间的 Quillen
\trEquivalence{}，通过明智地选取单纯富集，可以理解为 $(\infty,2)$-\trCategory——而不仅是
$(\infty,1)$-\trCategory——的\trEquivalence{}（见 E.2 节）。}由此可知，[78] 中以分析式手段
证明的关于\trQuasiCategory{}的定理，对\trCompleteSegalSpaces{}同样成立，反之亦然。我们以这一
迁移原理（transfer principle）的若干应用收尾：例如，在对象为 $(\infty,1)$-
\trCategory{}的\trInfinityCosmos{}中，我们证明各种泛性质都是「逐点确定」的——先用分析式技术对%
\trQuasiCategory{}证明这些结果，再诉诸\trModelIndependence{}，把它们\trExtension{}到\trBiequivalence{}的\trInfinityCosmos{}。

关于\trInfinityCategory{}的陈述之模型不变性的问题，比人们预想的更为微妙。把一个\trInfinityCategory{}从
一个模型搬到另一个模型再搬回来时，所得对象通常与原物\trEquivalence{}而不恒同；\trInfinityCategory
的某些「邪恶」（evil）性质在\trEquivalence{}之下并非不变——「一个\trInfinityCategory
只有单个对象」这一断言即是著名的例子。\yizhu{「邪恶性质」（evil
property）是范畴论行话，指在\trEquivalence{}之下不保持、因而不堪当范畴论性质之名的
属性。}我们这套系统理解\trInfinityCategory{}理论之\trModelIndependence{}的取径，其一大优点是允许我们
引入一种形式语言（formal language），并证明在该语言中可表达的关于\trInfinityCategory{}的
论断都是模型无关的。这建立在 Makkai 的工作之上——他解决了一个类似的问题，
即 2-\trCategory{}的性质在\trBiequivalence{}之下的不变性 [82]。

遗憾的是，限于篇幅，我们未能探访\trInfinityCosmos{}中那些\trHomotopy{}相干的结构。例如，一篇
姊妹论文 [109] 证明：\trInfinityCosmos{}中\trInfinityCategory{}之间的任何\trAdjunction{}都\trHomotopy{}唯一地%
\trExtension{}成一个\trHomotopy{}相干\trAdjunction{}，并给出了\trHomotopy{}相干单子（monad）的单子性
定理，作为\trInfinityCategory{}式泛代数（universal algebra）的一种机制。\trHomotopy{}相干单子的
形式理论又由 Sulyok [124] 进一步拓展——他发展了相应的单子下降与余单子
下降（descent）理论；Zaganidis [133] 亦有贡献——他定义并研究了\trHomotopy{}相干
单子映射。篇幅所限的另一位牺牲者，是对\trCartesianFibration{}的「大宇宙
原理」（macrocosm principle）的探索：该原理断言，从\trCartesianFibration{}的\trInfinityCosmos
指向底\trInfinityCosmos{}的陪域投射\trFunctor{}，在某个适当的意义下定义了一个「\trInfinityCosmos{}的%
\trCartesianFibration」[111]。我们希望在续篇中回到这些课题。

理想的读者或许已对富集\trEnrichedCategory{}论、2-范畴论与抽象\trHomotopy{}论有所
涉猎，于是那些在这些传统中自有渊源的构造与证明读来便会感到熟悉。由于%
\trInfinityCategory{}吸引着背景极为多样的数学家，我们在附录 A、B 与 C 中分别温习了这些
主题所需的全部材料。要在第 1 章中确立一般\trInfinityCosmos{}的相应特征，需要 Joyal
最先证明的关于\trQuasiCategory{}的若干基本事实。我们在 1.1 节中陈述这些结果，而把那些
需要冗长组合枝蔓的证明推迟到附录 D——在那里我们还引入了 \emph{n}-复状集，
它对任意 $0\le n\le\infty$ 给出 $(\infty,n)$-\trCategory{}的一个模型。「在野外」
出现的\trInfinityCosmos{}实例可在附录 E 中找到；那里还介绍了读者可用以自己发掘\trInfinityCosmos
的一般技术。最后一个附录则处理一桩关键的未竟之事：至关重要的一点是，在拟%
\trCategory{}的\trInfinityCosmos{}中发展出的综合式理论，与 Joyal、Lurie 以及众多学者发展的
分析式理论完全相容。这是附录 F 的主题。

我们以关于大小的例行免责声明作结。要把此处发展的理论应用到最有意思的那些%
\trInfinityCategory{}上，就应当考虑三个无穷不可达基数（inaccessible cardinals）
$\alpha<\beta<\gamma$，一如通行的约定 [5, 2]。通俗地说，$\alpha$-小的\trCategory
或可称为「小」的，而 $\beta$-小的\trCategory{}则是\trInfinityCategory{}的默认尺寸。例如，由
（小）空间、（小）阿贝尔群（abelian groups）的链复形（chain complexes）或
（小）\trHomotopy{}相干图构成的\trInfinityCategory{}都是 $\beta$-小的。这些常规尺寸的\trInfinityCategory{}，正是
一个 $\gamma$-小的\trInfinityCosmos{}的对象——用通俗的话说，「大」的。当然，倘若只对小%
\trSimplicialSet{}感兴趣，那么小\trQuasiCategory{}构成的\trInfinityCosmos{}便是 $\beta$-小的而非
$\gamma$-小的，此处发展的理论对它照样适用。因此，我们把 Grothendieck 宇宙
（Grothendieck universes）搁置一旁，此后不再提及这些不可达基数。
